% Part: conditional-logics
% Chapter: minimal-change-semantics
% Section: true-false

\documentclass[../../../include/open-logic-section]{subfiles}

\begin{document}

\olfileid{con}{min}{tf}

\olsection{反事实条件句的真与假}

一个反事实条件句 $!A \cif !B$ 为（非空）真，如果最近的 $!A$-世界都是 $!B$-世界，如 \olref{fig:true} 所示。

\begin{figure}
\begin{center}
\begin{tikzpicture}[scale=.7]\small
  \spheresystem{5}
  \draw (0,0) node {$w$};
  \propositionintersect{0}{4}{40}{3.5}
  {\tikzset{proposition/.append={smooth,tension=1.4}}
  \proposition[shift={(1.6,-1)}]{90}{1}{30}{4}}
  \path (1.5,3) node[below] {$\formula{B}$};
  \path (3,0) node {$\formula{A}$};
\end{tikzpicture}
\caption{非空真的反事实条件句}
\ollabel{fig:true}
\end{center}
\end{figure}

如果环绕 $w$ 的球系中根本没有任何容纳 $!A$ 的球，则反事实条件句在 $w$ 处也为真。此时，它是\emph{空真}的（见 \olref{fig:vacuous}）。

\begin{figure}
\begin{center}
\begin{tikzpicture}[scale=.7]\small
  \spheresystem{5}
  \draw (0,0) node {$w$};
  \proposition{0}{6.5}{40}{3.5}
  {\tikzset{proposition/.append={smooth,tension=1.4}}
  \proposition[shift={(1.2,-1)}]{90}{1}{30}{4}}
  \path (1.5,3) node[below] {$\formula{B}$};
  \spherepos{0}{7}{node {$\formula{A}$}}
\end{tikzpicture}
\caption{空真的反事实条件句}
\ollabel{fig:vacuous}
\end{center}
\end{figure}

反事实条件句为假有两种情况。其一，最近的 $!A$-世界并非都是 $!B$-世界，但其中有些是。此时，$!A \cif \lnot !B$ 也为假（见 \olref{fig:false}）。

\begin{figure}
\begin{center}
\begin{tikzpicture}[scale=.7]\small
  \spheresystem{5}
  \draw (0,0) node {$w$};
  \propositionintersect{0}{4}{40}{3.5}
  {\tikzset{proposition/.append={smooth,tension=1.4}}
  \proposition[shift={(1.6,0)}]{90}{1}{30}{3}}
  \path (1.5,3) node[below] {$\formula{B}$};
  \path (3,0) node {$\formula{A}$};
\end{tikzpicture}
\caption{假的反事实条件句，对立句亦假}
\ollabel{fig:false}
\end{center}
\end{figure}

如果最近的 $!A$-世界与 $!B$-世界完全不重叠，那么 $!A \cif !B$ 为假。但在这种情况下，所有最近的 $!A$-世界都是 $\lnot !B$-世界，因此 $!A \cif \lnot !B$ 为真（见 \olref{fig:false-opposite}）。

\begin{figure}
\begin{center}
\begin{tikzpicture}[scale=.7]\small
  \spheresystem{5}
  \draw (0,0) node {$w$};
  \propositionintersect{0}{4}{40}{3.5}
  {\tikzset{proposition/.append={smooth,tension=1.4}}
  \proposition[shift={(-1,0)}]{90}{1}{30}{3}}
  \path (-1,3) node[below] {$\formula{B}$};
  \path (3,0) node {$\formula{A}$};
\end{tikzpicture}
\caption{假的反事实条件句，对立句为真}
\ollabel{fig:false-opposite}
\end{center}
\end{figure}

与严格条件句不同，反事实条件句可以是偶真的。考虑 \olref{fig:contingent} 中的球系模型。距离 $u$ 最近的 $!A$-世界都是 $!B$-世界，所以$\mSat{M}{!A \cif
  !B}[u]$。但是，距离 $v$ 最近的某些 $!A$-世界不是 $!B$-世界，所以$\mSat/{M}{!A \cif !B}[v]$。

\begin{figure}
\begin{center}
\begin{tikzpicture}[scale=.6]\small
  \spheresystem{7}
  \spheresystem[shift={(0:3.1)},dashed]{4}
  \propositionintersect{45}{4}{40}{4.5}
  \begin{scope}
    \clip \propositionplot{45}{4}{40}{4.5} ;
    \spherelayer[shift={(0:3.1)}]{3}
  \end{scope}
  \proposition[shift={(1.4,-.5)}]{90}{1}{35}{5}
  \path (0,0) node {$u$};
  \path (0:3.1) node {$v$};
  \spherepos{35}{9}{node {$\formula{A}$}}
  \spherepos{80}{9}{node {$\formula{B}$}}
\end{tikzpicture}
\end{center}
\caption{偶真的反事实条件句}
\ollabel{fig:contingent}
\end{figure}

\end{document}